\subsubsection{Razonamiento Cronológico y Extracción de Datos en Conversaciones Familiares}

En la notebook 11 se implementa un experimento para extraer y organizar datos de conversaciones con un grupo familiar, buscando corregir el razonamiento cronológico del modelo de lenguaje GPT-3.5. A continuación se detallan los pasos principales del experimento.

\paragraph{Carga de Librerías y Configuración}

Se utiliza FAISS para realizar la búsqueda semántica y OpenAI Embeddings para convertir los textos de las conversaciones en vectores. También se configura una función para interactuar con el modelo GPT-3.5 a través de la API de OpenAI.

\begin{APACode}
import os
import requests
import json
from langchain.embeddings.openai import OpenAIEmbeddings
from langchain.vectorstores import FAISS

OPENAI_API_KEY = os.getenv("OPENAI_API_KEY")

index_name = "conversations"
index_path = "../indexes/conversations"
embeddings = OpenAIEmbeddings()

def gpt_system_user(
    system_message: str, user_message: str, model: str = "gpt-3.5-turbo"
):
    """Para usar en notebooks"""
    try:
        response = requests.post(
            url="https://api.openai.com/v1/chat/completions",
            params={
                "temperature": "0",
            },
            headers={
                "Content-Type": "application/json",
                "Authorization": f"Bearer {OPENAI_API_KEY}",
            },
            data=json.dumps(
                {
                    "model": model,
                    "messages": [
                        {"content": system_message, "role": "system"},
                        {"content": user_message, "role": "user"},
                    ]
                }
            ),
            timeout=10,
        )
        choices = response.json()["choices"]
        answer = {
            "status_code": response.status_code,
            "text": choices[0]["message"],
        }
        return answer
    except requests.exceptions.RequestException:
        print("HTTP Request failed")
        return None
\end{APACode}

\paragraph{Simulación de Conversaciones Familiares}

Se crean varias conversaciones simuladas en las que el usuario proporciona información sobre su familia, como el hecho de tener una hija mayor llamada Ana de 20 años, y un hijo menor llamado Juan de 10 años. También se incluyen detalles de un viaje familiar a Córdoba con fechas y otros aspectos relevantes.

Ejemplos de mensajes:
\begin{itemize}
    \item Usuario: ``Voy a sacar un pasaje para la familia de Buenos Aires a Córdoba''.
    \item Asistente: ``¿Quiénes son los pasajeros?''.
    \item Usuario: ``Ana es mi hija mayor, tiene 20 años. Juan es mi hijo menor, tiene 10 años''.
\end{itemize}

\begin{APACode}
condensed_contexts = [
    {
        "messages": [
            {
                "role": "user",
                "content": "Voy a sacar un pasaje para la familia de buenos aires a cordoba"
            },
            {
                "role": "system",
                "content": "quienes son los pasajeros?"
            },
            {
                "role": "user",
                "content": "Ana es mi hija mayor, tiene 20 años. Juan es mi hijo menor, tiene 10 años. Mi esposa se llama María"
            },
            {
                "role": "system",
                "content": "en que fecha quiere viajar?"
            },
            {
                "role": "user",
                "content": "del 20 de agosto al 30 de agosto"
            },
            {
                "role": "system",
                "content": "perfecto, salen el 20 de agosto a las 10:00 y vuelven el 30 de agosto a las 18:00. El precio total es $500"
            },
            {
                "role": "user",
                "content": "si"
            },
            {
                "role": "system",
                "content": "Gracias por su compra"
            }
        ],
        "metadata": {
            "session_id": "1",
            "person_id": "20",
            "context_id": "1",
            "created_at": "20200825T12:00:00.000Z",
        },
    }
]
\end{APACode}

\paragraph{Extracción de Datos}

Se utiliza un prompt diseñado para que el modelo GPT-3.5 extraiga la información relevante de cada conversación, incluyendo nombres, edades y fechas de eventos. Esta información se almacena junto con los metadatos en FAISS para facilitar su búsqueda posterior.

\begin{APACode}
db = FAISS.from_texts(texts=['notebook-11'], embedding=embeddings)
for context in condensed_contexts:
    messages = ''.join([str(message) for message in context['messages']])
    extraction = gpt_system_user(
        system_message="""
Extract every meaningful information about the user from the provided conversation.
Ask yourself, what happened in the conversation?
What data about the user did you learn?
At what date and time was this data gathered? Use this response format:
{
    DATA: `the meaningful information that you have learned about the user`,
    CREATED_AT: `the date and time when this data was gathered`,
    RELATIONSHIP: `the relationship between the user and the data that you have learned`,
}

Take a deep breath and work on this problem step-by-step.
        """,
        user_message=f"METADATA: ```{context['metadata']}``` CONVERSATION: ```{messages}```",
        model="gpt-3.5-turbo",
    )
    db.add_texts(
        texts=[extraction['text']['content']],
        metadatas=[context['metadata']],
    )
\end{APACode}

Ejemplo de extracción:
``Voy a sacar un pasaje para la familia de Buenos Aires a Córdoba, Ana es mi hija mayor, tiene 20 años. Juan es mi hijo menor, tiene 10 años.''

\paragraph{Consultas sobre la Edad}

Se realiza una consulta al sistema preguntando ``¿Qué edad tiene mi hija mayor?''. El sistema recupera la información y responde correctamente que Ana tiene 20 años. Sin embargo, no razona correctamente sobre el año actual, que es 2023, lo que significa que Ana debería tener 23 años. Este error revela una falta de ajuste cronológico en el razonamiento del modelo.

\begin{APACode}
user_input_3 = "que edad tiene mi hija mayor?"
filter = {"person_id": "20"}
db_results_3 = db.similarity_search_with_score(
    query=user_input_3, embeddings=embeddings, filter=filter, k=1
)
context_3 = "{ message: " + db_results_3[0][0].page_content + ", metadata: " + str(db_results_3[0][0].metadata) + " }"

system_prompt_3 = f"""
You are a personal assistant. You will find past conversations between you and the user between backticks.
Use them to answer the user's question.
To reason about dates, have in mind that today is September 13th, 2023.
You are now located in Buenos Aires, Argentina.
Take a deep breath and work on this problem step-by-step.
```{context_3}```
"""

answer_3 = gpt_system_user(system_message=system_prompt_3, user_message=user_input_3)
print(answer_3["text"]["content"])
\end{APACode}

\paragraph{Conclusiones de la undécima iteración}

A pesar de proveer información contextual sobre la fecha actual en los prompts, el sistema sigue cometiendo errores al no actualizar adecuadamente la edad de Ana basándose en el año actual. El modelo responde incorrectamente con la información de 2020, sin hacer el cálculo necesario para ajustarla a la fecha actual.

La respuesta del sistema fue: ``Tu hija mayor tiene 20 años.''

Persiste el problema de razonamiento sobre la fecha. El dato de la edad es de 2020, la respuesta correcta es 23 años considerando que la consulta se realiza en 2023.

La notebook 11 muestra que, si bien el modelo GPT-3.5 puede extraer correctamente datos de conversaciones y responder a preguntas basadas en esa información, sigue teniendo dificultades para razonar cronológicamente. Aunque se le proporciona contexto sobre las fechas actuales, el modelo no ajusta correctamente la edad de las personas en función del tiempo transcurrido. Esto indica la necesidad de mejorar los mecanismos de razonamiento temporal en los modelos de lenguaje.

\textit{Código disponible en:} \url{https://github.com/aditzend/hyperpersonalization/blob/main/app/notebooks/11.ipynb} 